\section{Reference model and equilibrium inputs}\label{sec:inputs}

\subsection{Reference conventions}\label{model:reference}

All conditional averaging below uses the \emph{critical nearest-neighbor}
measure at bond coupling $K_c$, in the plane or with the indicated physical
boundary condition. We write $\mu^0$ for such a reference law and identify
which geometry is meant whenever it matters. For a finite cut, the
conditional law is determined by its cutting spins and physical boundary
data; it is not a conditional law of the interacting system.

After rescaling a local chart, its fine mesh is denoted by $\delta$.
For the renormalization construction we return to lattice-unit lengths.
There $B_r(z)$ is a max-norm square centered at $z$, and $P_{r,z}F$ is
the conditional expectation of $F$ given the exterior and the cutting
contour of $B_r(z)$. A consistent rounding convention includes the cutting
vertices in the retained exterior. The resulting function depends only on
those vertices, together with the prescribed physical boundary. Nested
cuts compose by conditional expectation. Bounded changes in lattice
rounding are absorbed in fixed buffer margins. Write $P_r=P_{r,0}$.

The microscopic scalar energy used for the distinguished direction is
\begin{equation}\label{model:scalar-energy}
 e(z)=\frac14\sum_{v:\,|v-z|=1}\sigma_z\sigma_v.
\end{equation}
Each unoriented bond occurs twice in this sum over $z$, so
$\sum_z e(z)$ is half the total bond-product sum. The source fields are
$\mathcal O_1=e$ and $\mathcal O_{1/8}=\sigma$, with the subscript
indicating scaling dimension. Individual directed bonds are retained in
the target; their common normalization will follow from square symmetry,
not from changing the observable in \eqref{model:energy}.

Supremum norms of local functions are actual suprema over their available
spins. Background interactions are considered modulo constants, represented
by subtracting the mean under independent uniform available spins. This
projection costs at most a factor two in the supremum norm. Source scalar
terms will be recorded separately. All constants may depend on fixed
buffer geometry. Once the scale ratio is chosen, analytic constants may
also depend on that ratio and the fixed potential; constants for source
jets may depend on the finite number of marks.

\subsection{Equilibrium inputs}

Every measure in this section is ferromagnetic. Except in the
finite-graph statement below, its nearest-neighbor coupling is $K_c$.
External fields introduced here are auxiliary fields in this reference
model. None of the comparison inequalities is asserted for the
perturbed Hamiltonian of Theorem~\ref{model:main}.

We state the particular consequences of the established reference
theory that will be used. The pointwise, ordered-transfer, and stopping
statements are proved in the companion article
\emph{Buffered comparison and stopping-band resolution in critical Ising}
\cite{IsingToolkit}. We state every interface used here, including its
restrictions on common pins and unchanged continuations. The
pointwise comparison uses the external-field inequality of
Ding--Song--Sun \cite[Corollary~1.3]{DSS2023}.

\begin{lemma}[Pointwise comparison with common pins]\label{in:ratio}
Let $G$ be a finite graph with nonnegative pair couplings and arbitrary
finite real fields. Let $I,J,F$ be disjoint vertex sets, with $I,J$
nonempty, and prescribe arbitrary common spin values on $F$. Form the
graph obtained by deleting $F$ and contracting $I$ and $J$ separately;
parallel couplings are added. Let $q_0$ be the probability that the two
contracted vertices are FK-connected in the zero-field, cluster-weight
$2$ measure on this graph. For any configurations $a$ on $I$ and $b,b'$
on $J$,
\begin{equation}\label{in:ratio-bound}
 e^{-4\operatorname{arctanh}q_0}
 \leq
 \frac{\Prob(\sigma_I=a\mid\sigma_J=b,\sigma_F)}
      {\Prob(\sigma_I=a\mid\sigma_J=b',\sigma_F)}
 \leq e^{4\operatorname{arctanh}q_0}.
\end{equation}
Either conditional law may be replaced by a mixture of such laws.
All external fields, including those produced by the common pins,
are set to zero in the comparator defining $q_0$.
\end{lemma}

This is \cite[Lemma~2.2]{IsingToolkit}. Common pins are deleted vertices in
the comparator; they are not additional FK wires. This distinction
will allow the same estimate in the presence of many signed pins.

\begin{corollary}[Buffered and shrinking-patch comparison]\label{in:patch}
Fix $c>0$. If $I\subset B_R(z)$ and
$J\subset B_{(1+c)R}(z)^c$, their conditional density ratios in
Lemma~\ref{in:ratio} lie in $[C_c^{-1},C_c]$, uniformly in the mesh,
common pins, and fields. There are $C<\infty$ and $\theta>0$ such that,
if $I\subset B_w(z)$, $J\subset B_R(z)^c$, and $R\geq8w\geq8$ in
lattice units, the absolute log density ratio is at most
\begin{equation}\label{in:patch-bound}
 C(w/R)^\theta.
\end{equation}
No unpinned neighborhood of $I$ is required.

The buffered bound holds for fixed polygonal separating collars with
positive relative width, and for their intersections with a free or
fixed physical boundary. It also compares different continuations or
volumes when the reference interaction agrees in the separating
chart. The same assertions hold for mixtures and after a common
conditional randomization of the local spins. Finitely many shrinking
patches can be changed successively.
\end{corollary}

\begin{proof}
For square collars this is
\cite[Corollaries~2.5--2.6]{IsingToolkit}.
Its proof uses the critical FK-Ising RSW estimate with arbitrary
boundary wiring \cite{DHN2011}. A finite collection of rectangular
dual crossings supplies a separating circuit with probability bounded
below. Repeated disjoint annular tests give
$q_0\leq C(w/R)^\theta$. In a clipped chart restore the missing
lattice vertices and edges when bounding $q_0$ from above, and then
apply the same tests in the restored collar. Common fixed boundary
spins have already been deleted. Polygonal collars use finitely many
such rectangular tests. Conditioning on a farther contour represents
each continuation as a mixture of the same local kernels. These
observations give the stated extensions of the square version.
\end{proof}

For clarity, we give the geometry needed in the ordered-transfer
statement. A cut $J$ separates a source $S_0$ from a target $T$ in the
graph of unpinned sites. At lattice scales $1\ll w\leq s$, call the
cut \emph{admissible} if it is partitioned into patches $b_i$, with
straight counting segments $L_i\subset J$ and centers $z_i$, such that
the following hold with fixed positive comparison constants. Each
$b_i\cup L_i$ lies in a square of radius $\ell_i\asymp w$;
$|L_i|\asymp\ell_i$; the square $B_{16\ell_i}(z_i)$ is an unpinned
lattice chart disjoint from the source; the $L_i$ have bounded overlap;
and $T$ is a distance comparable to $s$ away, with
$16\ell_i$ smaller than that distance. The patch estimate must be
available up to a fixed multiple of $s$. A square contour in a
pin-free collar of width comparable to $w$ satisfies these conditions
after subdivision, including at its corners. Put
\[
 M_J=s^{-7/8}\sum_{x\in J}\sigma_x.
\]

\begin{corollary}[Ordered changes with a common continuation]
\label{in:ordered}
Suppose $\mu^-\preceq\mu^+$ are two ordered laws, possibly on
different volumes or with different pins. Assume there is a common
unpinned source contour $S_0$ such that, given its spins, both laws
have the same critical nearest-neighbor continuation toward $T$.
Apart from the common pinning data, this continuation has zero field.
All changed data lie on the source side of $S_0$. For an admissible
intermediate cut $J$, there are $C<\infty$ and
$\theta\in(0,1/8)$ such that
\begin{equation}\label{in:ordered-bound}
 \norm{\mu_T^+-\mu_T^-}_{\TV}
 \leq C(w/s)^{-7/8+\theta}
       \bigl(\E_{\mu^+}M_J-\E_{\mu^-}M_J\bigr).
\end{equation}
Here total variation is the supremum of event-probability differences.
The constants depend on the fixed chart and admissibility constants,
not on the source values or their number. At a fixed positive $w/s$,
vanishing rescaled magnetic difference implies vanishing target total
variation. Random ordered source configurations are allowed.
\end{corollary}

This is \cite[Corollary~4.12, with Definition~4.1 and
Corollary~4.11]{IsingToolkit}.
Fields or additional changes before $S_0$ are allowed by first coupling
the induced source distributions by order and then using the common
continuation. The target-side continuation has no auxiliary field,
and its pinning data are identical in the two laws.

\begin{theorem}[Outward stopping-band transfer]\label{in:transfer}
Work at mesh $\delta$ in a fixed square chart with homogeneous fixed
sign on its outer wall. Fix nested squares $Q_u\Subset Q_v$ and a
target $T$ outside $Q_v$ and inside the wall, with positive clearances
from both. Reserve a slightly larger band window, also separated from
$T$ and the wall. The levels $u,v$ may be chosen generically in
prescribed separated subintervals of this window.

Initially reveal $Q_u$ and its cutting contour. There is a finite
adaptive spin exploration through the band, with full transcript
$\mathcal F^\delta$, such that the unexamined sites, conditional on
each realizable transcript, have the ordinary Gibbs law with the
examined spins fixed. Given $\epsilon,\eta>0$, this observation has a
descriptor
\[
 Z^\delta\in\{0,1,\ldots,m\}
\]
computed from band spins, with the following properties for all
sufficiently small $\delta$:
\begin{enumerate}
\item $\Prob(Z^\delta=0)\leq\epsilon$ in the fixed-wall reference law.
\item For every $i\geq1$ there is a deterministic regular signed
Jordan barrier $\Gamma_i$, with finitely many sign arcs, on the fresh
side of the target-facing revealed boundary. It separates the reveal
from $T$ and has a positive clearance from $T$ and the distant wall.
For \emph{every} realizable transcript with $Z^\delta=i$, adding the
plus barrier pins and then the minus barrier pins changes the target
law by a sum of at most $\eta$ in total variation.
\item After these additions, the target-side law is a single law
$Q_{i,\delta}$, determined by $\Gamma_i$ and the unchanged outer
wall, independent of the particular transcript in bin $i$.
\end{enumerate}
The finite list, all its positive clearances, and $m$ are chosen
before the mesh threshold, depending only on the geometry and
$\epsilon,\eta$. They work simultaneously for all bounded target
tests, including mesh-dependent tests. For a finite list of bands,
the descriptors may be chosen with converging joint laws.
If the band-spin distribution under another law has density at most
$C$ relative to the reference distribution, its exceptional
probability is at most $C\epsilon$.
\end{theorem}

The statement is the outward part of
\cite[Theorem~9.1 and Lemma~7.2]{IsingToolkit}. In particular, for
$|H_\delta|\leq1$, on every good transcript,
\begin{equation}\label{in:lookup}
 \left|\E[H_\delta\mid\mathcal F^\delta]
        -\E_{Q_{Z^\delta,\delta}}H_\delta\right|\leq2\eta.
\end{equation}
The barriers can be polygonal or regular smooth contours. They need
not lie between the two numerical levels of an individual band;
the reserved window accommodates their offsets and tip caps.
The last assertion transfers the \emph{exceptional probability}.
A change of target-side continuation requires an additional buffered
comparison; this is carried out in Section~\ref{sec:buffers}.

\begin{proposition}[Static critical-field inputs]\label{in:fields}
The following facts concern only the critical nearest-neighbor model.
\begin{enumerate}
\item\label{in:fields-convergence}
For bounded finitely connected domains without singleton boundary
components and with finitely many nondegenerate fixed-plus, fixed-minus,
or free boundary arcs, separated bulk spin and
energy correlations converge with dimensions $1/8$ and $1$. Discrete
domains converge in the Caratheodory sense, with eventually matching
boundary-component and arc counts and converging boundary marks.
Convergence is uniform on compact separated configurations. The limits
are conformally covariant and continuous under these marked-domain
limits. The energy convention initially subtracts its plane mean;
expanding the finitely many subtractions gives the same assertion for
domain-centered energies. The two bond directions have the same
scalar limit. Homogeneous half-plane spin limits and full-plane spin
limits are included, with odd plane correlations zero.
\item\label{in:fields-formulas}
Normalize the continuum spin so that its plane two-point function is
$|z-w|^{-1/4}$. Write $G_n$ for its $n$-point function. For even $n$,
\begin{equation}\label{in:plane-spin}
 G_{n,\C}(z_1,\ldots,z_n)^2
 =2^{-n/2}
   \sum_{\substack{\mu\in\{-1,1\}^n\\\sum_i\mu_i=0}}
       \prod_{i<j}|z_i-z_j|^{\mu_i\mu_j/2}.
\end{equation}
For $z_i$ in the upper half-plane, put
$q_{ij}=|(z_i-z_j)/(z_i-\overline z_j)|$. Then
\begin{align}
 G_{n,\mathbb H}^{+}(z)^2
 &=2^{-n/2}\prod_i(2\operatorname{Im}z_i)^{-1/4}
       \sum_{\mu\in\{-1,1\}^n}\prod_{i<j}q_{ij}^{\mu_i\mu_j/2},
       \label{in:halfplane-plus}\\
 G_{n,\mathbb H}^{\mathrm{free}}(z)^2
 &=c_n\prod_i(2\operatorname{Im}z_i)^{-1/4}
       \sum_{\mu\in\{-1,1\}^n}
          \left(\prod_i\mu_i\right)
          \prod_{i<j}q_{ij}^{\mu_i\mu_j/2},\quad n\ \hbox{even}.
       \label{in:halfplane-free}
\end{align}
Here $c_n=(-1)^{n/2}2^{-n/2}$ for even $n$, including $c_0=1$.
The plus formula holds also for odd $n$; minus boundary
is obtained by spin flip. These functions inherit reflection
positivity from the lattice and satisfy spherical covariance in the
plane. The plus-disk one-point functions of the spin and the
plane-centered scalar energy are nonzero.
\item\label{in:fields-moments}
For homogeneous plus, minus, or free reference laws and for the plane,
the rescaled two-point bound is
\[
 \delta^{-1/4}|\E\sigma_z\sigma_{z'}|
 \leq C(|z-z'|\vee\delta)^{-1/4}
\]
on buffered compact sets. More generally, if $z_1,\ldots,z_n$ lie in
a fixed compact interior set and
$\rho_i$ is the lesser of a fixed buffer radius and one quarter of
the distance to the other insertions, then
\begin{equation}\label{in:product-arm}
 \left|\E\prod_i\sigma_{z_i}\right|
 \leq C_n\prod_i
       \left(\frac{\delta}{\rho_i\vee\delta}\right)^{1/8}.
\end{equation}
Repeated sites are allowed with the factor capped at one. The estimate
uses ordinary FK wiring and homogeneous cluster coloring; prescribed
mixed boundary signs are handled by buffered density comparison.
For bounded compactly supported piecewise regular profiles $h$, the
variables
$X_\delta(h)=\delta^{15/8}\sum_z h(z)\sigma_z$
have uniform exponential moments of every fixed real order under
homogeneous reference laws on buffered compact sets. Positive-profile
bounds follow from GHS; signed profiles follow by splitting the signs
and applying Cauchy--Schwarz. These moment inputs also hold in
homogeneous half-plane references by local comparison.
\item\label{in:fields-necklace}
In a fixed square with free boundary and zero field, for every
$\epsilon>0$ there are $K<\infty$ and $c>0$ such that, for all
sufficiently small $\delta$, with probability at least $1-\epsilon$
a nearest-neighbor vertex circuit surrounds the concentric square of
one third the side length. The circuit is contained in the union of
at most $K$ FK clusters, each of cardinality at least
$c\delta^{-15/8}$; successive clusters may be separated by one
nearest-neighbor step. For a uniform positive field
$h\delta^{15/8}$, conditional on internal FK bonds, a cluster $C$
is attached to the plus ghost with probability
\begin{equation}\label{in:ghost}
 \tanh(h\delta^{15/8}|C|),
\end{equation}
independently for distinct internal clusters. Internal bonds dominate
their zero-field law. Consequently, at fixed physical square size,
the probability of a plus-spin circuit tends to one by choosing a
sufficiently large finite $h$ and then a sufficiently small mesh.
\end{enumerate}
\end{proposition}

The bounded-domain convergence is
\cite[Theorem~1.3 and Definitions~3.1--3.2]{CHI2022}; its
Theorem~7.1 gives \eqref{in:halfplane-plus}. The plane neutral-charge
form is \cite[Theorem~4]{Dubedat2011}, with the displayed normalization
agreeing with the half-plane formula as the boundary recedes.
The homogeneous free formula also follows directly from
\cite[Theorem~1.1 and Definition~2.3]{BIVW2025}: the constant boundary
shift $\pi/\sqrt2$ changes each bosonic cosine into a negative sine,
whose charge expansion has coefficient $i\mu_i/\sqrt2$ for charge
$\mu_i\in\{-1,1\}$. Multiplication over an even number of insertions
gives $c_n\prod_i\mu_i$ in \eqref{in:halfplane-free}. The half-plane
Green function $-\log q_{ij}$ and its regularized diagonal
$\log(2\operatorname{Im}z_i)$ supply the displayed pair and height
factors. The bounded-domain
theorem and exhausting homogeneous domains supply half-plane
references; changing a distant wall is controlled by
Corollary~\ref{in:patch}. The nonzero plane-centered energy response in a fixed-plus disk is
also explicit in \cite[Theorem~1.1 and Corollary~1.2]{HS2013}.
No rate of primary-field convergence is used.

Continuity under the stated continuum marked-domain limits follows
from the same discrete convergence theorem. Given a nondegenerate
sequence $D_j\to D$, choose lattice approximations of $D_j$ at meshes
$\delta_j<1/j$ for which the rescaled correlations differ from their
continuum values by less than $1/j$. Choose them sufficiently close
to $D_j$ that the diagonal sequence converges to $D$ with the required
component and arc counts and boundary marks. Applying the discrete
convergence theorem to this sequence proves continuity. Allowing the
marked insertion configurations to vary in a compact separated set
gives its locally uniform version. This deduction keeps the prescribed
topology and nondegenerate limiting boundary arcs.

For the third item, the FK one-arm estimate and product bound are
recorded in \cite[Lemmas~3.1--3.3]{IsingToolkit}. The exponential-moment input is
\cite[Section~3.1, Proposition~3.5, Theorem~3.6,
Corollaries~3.7--3.8]{CGN2015}; the passage from these inputs to the
integrated limits used here is proved in Lemma~\ref{buf:moments}.
The last item is \cite[arXiv version~4, Lemmas~4 and~8 and the proof of
Lemma~9]{CJN2020}. Indeed its success probability is bounded below by
$(1-\epsilon)\tanh(ch)^K$, after an arbitrary initial reduction of
$\epsilon$. The event may be formulated using connected subsets of
the required size, so merging clusters preserves the geometric event.

We use finite-volume positive association and stochastic ordering
\cite{FKG1971,Holley1974}, the GHS inequality \cite{GHS1970}, the
Edwards--Sokal coupling \cite{EdwardsSokal1988}, and Euclidean reflection
positivity \cite{FILS1978} only under the hypotheses stated above. Whenever an infinite-volume reference is used, the finite
comparison is first applied inside a chart and then its outer wall
is sent to infinity. All its constants are independent of that wall.
